% Source: Weinberg GR, physical PDF pp. 6--8
%         (printed Preface pp. vii--ix).
% Coverage: Preface.
% Figures/tables/footnotes: none.
% Status: source-reviewed and compile-clean.
% Uncertainties: none.

\clearpage
\section*{Preface}
\addcontentsline{toc}{section}{Preface}

Now that this book is done I can look back, and identify two purposes which
led me to begin writing, and which have guided the work to completion.

One good practical purpose was to bring together and assess the wealth of data
provided over the last decade by new techniques in experimental physics and in
optical, radio, radar, X-ray, and infrared astronomy. Of course, new data will
keep coming in even as the book is being printed, and I cannot hope that this
work will remain up to date forever. I do hope, however, that by giving a
comprehensive picture of the experimental tests of general relativity and
observational cosmology, I will help to prepare the reader (and myself) to
understand the new data as they emerge. I have also tried to look a little way
into the future, and to discuss what may be the next generation of experiments,
especially those based on artificial satellites of the earth and sun.

There was another, more personal reason for my writing this book. In learning
general relativity, and then in teaching it to classes at Berkeley and M.I.T., I
became dissatisfied with what seemed to be the usual approach to the subject. I
found that in most textbooks geometric ideas were given a starring role, so
that a student who asked why the gravitational field is represented by a metric
tensor, or why freely falling particles move on geodesics, or why the field
equations are generally covariant would come away with an impression that this
had something to do with the fact that space-time is a Riemannian manifold.

Of course, this was Einstein's point of view, and his preeminent genius
necessarily shapes our understanding of the theory he created. However, I
believe that the geometrical approach has driven a wedge between general
relativity and the theory of elementary particles. As long as it could be
hoped, as Einstein did hope, that matter would eventually be understood in
geometrical terms, it made sense to give Riemannian geometry a primary role in
describing the theory of gravitation. But now the passage of time has taught us
not to expect that the strong, weak, and electromagnetic interactions can be
understood in geometrical terms, and too great an emphasis on geometry can only
obscure the deep connections between gravitation and the rest of physics.

In place of Riemannian geometry I have based the discussion of general
relativity on a principle derived from experiment: the Principle of the
Equivalence of Gravitation and Inertia. It will be seen that geometric objects
such as the metric, the affine connection, and the curvature tensor naturally
find their way into a theory of gravitation based on the Principle of
Equivalence, and of course one winds up in the end with Einstein's general
theory of relativity. However, I have tried here to put off the introduction of
geometric concepts until they are needed, so that Riemannian geometry appears
only as a mathematical tool for the exploitation of the Principle of
Equivalence and not as a fundamental basis for the theory of gravitation.

This approach naturally leads us to ask \emph{why} gravitation should obey the
Principle of Equivalence. In my opinion the answer is not to be found in the
realm of classical physics, and certainly not in Riemannian geometry, but in
the constraints imposed by the quantum theory of gravitation. It seems to be
impossible to construct any Lorentz invariant quantum theory of particles of
mass zero and spin two unless the corresponding classical field theory obeys
the Principle of Equivalence. Thus the Principle of Equivalence appears as the
best bridge between the theories of gravitation and of elementary particles.
The quantum basis for the Principle of Equivalence is briefly touched upon here
in a section on the quantum theory of gravitation, but it was not possible to
go far into the quantum theory in this book.

The nongeometrical approach taken in this book has to some extent affected the
choice of the topics to be covered. In particular, I have not discussed in
detail \emph{the derivation and classification of complicated exact solutions
of the Einstein field equations}, because I did not feel that most of this
material was needed for a fundamental understanding of the theory of
gravitation, and hardly any of it seemed to be relevant to experiments that
might be carried out in the foreseeable future. By this omission I have left
out much of the work done by professional general relativists over the past
decade, but I have tried to provide an entree to this work through references
and bibliographies. I regret the omission here of a detailed discussion of the
beautiful theorems of Penrose and Hawking on gravitational collapse; these
theorems are briefly discussed in Sections 11.9 and 15.11, but an adequate
discussion would have taken up too much time and space.

I have tried to give a comprehensive set of references to the experimental
literature on general relativity and cosmology. I have also given references
to detailed theoretical calculations whenever I have quoted their results.
However, I have not tried to give complete references to all the theoretical
material discussed in the book. Much of this material is now classical, and to
search out the original references would be an exercise in the history of
science for which I did not feel equipped. The mere absence of literature
citations should not be interpreted as a claim that the work presented is
original, but some of it is.

It is a pleasure to acknowledge the inestimable help I have received in writing
this book. Students in my classes over the past seven years have by their
questions and comments helped to free the calculations of errors and
obscurities. I especially thank Jill Punsky for carefully checking many of the
derivations. I have drawn very heavily on the knowledge of many colleagues,
including Stanley Deser, Robert Dicke, George Field, Icko Iben, Jr., Arthur
Miller, Philip Morrison, Martin Rees, Leonard Schiff, Maarten Schmidt, Joseph
Weber, Rainier Weiss, and especially Irwin Shapiro. Finally, I am greatly
indebted to Connie Friedman and Lillian Horton for typing and retyping the
manuscript with inexhaustible skill and patience.

\begin{flushright}
\textsc{Steven Weinberg}
\end{flushright}

\noindent Cambridge, Massachusetts\\
April 1971
